\nocite{eda79af1}
\nocite{fd8cbb36}
\nocite{1706ae98}
The 2016 Nobel laureates in physics are Thouless, Haldane, and Kosterlitz {\glqq}for the theoretical discoveries of topological phase transitions and topological phases of matter{\grqq}, as officially stated by the Nobel Foundation.
Shortly after von Klitzing observed a quantized Hall conductance of a two-dimensional system of electrons at nearly zero temperature in a perpendicular magnetic field in 1980, the 2016 Nobel laureates showed that this very robust quantization does not merely originate in ordinary quantum mechanics, but rather in the topology of the spectrum.
\\
Von Klitzing found that the Hall conductance satisfies
\begin{align*}
  \sigma
  &=
  \frac{e^{2}}{h}
  \nu
  ,\qquad
  \nu
  \in
  \mathbb{N}
\end{align*}
where $e$ is the elementary charge and $h$ is Planck's constant.
$\nu$ is an integer depending monotonically on the strength of the magnetic field and is zero when the magnetic field vanishes.
Consequently, one calls the effect above the \textit{integer quantum Hall effect (IQHE)}.
It turned out that $\nu$ is a well-known topological invariant of the spectrum related to K-theory, making the IQHE the historically (compare \cite{1706ae98}) first example of a topological phase distinction.
If the IQHE is exhibited in a solid insulator, the insulator is called a topological insulator.
\\
The theory of topological insulators is a young field of research that utilizes many mathematical theories, which is probably the reason why there are not many convenient ways to engage with it - for the typical undergraduate student, at least.
The main purpose of this paper is to make up for this by serving as an accessible introduction to topological insulators.
\\
The paper is structured as follows:
\begin{enumerate}
\item[$\bullet$]
Section \ref{sec:mathpre} discusses the mathematical background needed to obtain a rough understanding of topological insulators.
The ultimate aim of Section \ref{sec:mathpre} is to equip the reader with topological K-theory which, in some limited cases, allows one to describe topological insulators.
\item[$\bullet$]
Section \ref{sec:topins} aims to convey a high-level picture of topological insulators.
We begin by stating precisely what we mean by an insulator.
After that, we briefly outline the general K-theoretic classification of topological phases.
The section culminates in the somewhat famous periodic table for topological insulators proposed in \cite{eda79af1} by Kitaev, as well as a high-level presentation of the IQHE and more general Chern insulators, which are contained in this periodic table.
\item[$\bullet$]
Section \ref{sec:kmm} - the last section - concerns the Kane-Mele model at a more detailed level.
It is a model for another kind of (two-dimensional) topological insulator contained in Kitaev's periodic table.
In a nutshell, we use methods similar to those of Section \ref{sec:topins} to describe the model and to find a topological invariant like the integer in the IQHE.
The treatment will be explicit enough to allow implementation, in contrast to the discussion in Section \ref{sec:topins}.
\end{enumerate}
